%! Representing a Quantitative Variable with Graphs — Unit 1, Lesson 1.3 — Warm-Up (Answer Key)
% ONE PAGE, blank and key. Three items, set at 12pt so they fill the page instead of leaving
% the bottom third empty. Do not add a fourth item — the page is full.
% Mirrors the blank exactly apart from the package swap, the pageheader, and the answers.
% NO teachernote here — teacher prose lives in the lesson plan.
\documentclass[12pt]{article}
\usepackage{apstats-article}
\usepackage{apstats-key}   % pulls in -boxes; do NOT also load -boxes

\begin{document}

\pageheader{Unit 1, Lesson 1.3}{Warm-Up --- Answer Key}

\vspace{0.28in}

\begin{enumerate}[leftmargin=1.9em, itemsep=0.80in]

  \item \textbf{From last lesson.} A theater sold 60 tickets on Friday and 240 on Saturday.
        On Friday 15 were student tickets; on Saturday 48 were.

        \vspace{12pt}
        Student share Friday: \ans{25\%} \qquad Saturday: \ans{20\%}
        \vspace{16pt}

        Which night had the larger \emph{share} of student tickets? \ans{Friday}

  \item Here are 12 numbers. Tally how many fall in each block.

        \vspace{12pt}
        \begin{center}
        21 \quad 34 \quad 27 \quad 45 \quad 38 \quad 22 \quad 49 \quad 31 \quad 26 \quad
        43 \quad 20 \quad 36
        \end{center}
        \vspace{12pt}

        \renewcommand{\arraystretch}{1.5}
        \begin{tabularx}{\linewidth}{@{} l Y Y Y Y @{}}
        \rowcolor{sky}
        \textbf{Block} & 20--29 & 30--39 & 40--49 & \textbf{Total} \\
        \hline
        \textbf{How many} & \ans{5} & \ans{4} & \ans{3} & \ans{12} \\
        \end{tabularx}

  \item Which of these could sensibly come out to \textbf{7.5}? \textbf{Circle} every one
        that could.
        \vspace{14pt}

        \hspace{1.0em}\textbf{(A)} the number of siblings a student has
        \hfill
        \ans{\textbf{(B)} a person's height in inches}
        \hspace{1.0em}
        \vspace{14pt}

        \hspace{1.0em}\textbf{(C)} the number of cars in a parking lot
        \hfill
        \ans{\textbf{(D)} the time to run one mile, in minutes}
        \hspace{1.0em}
        \vspace{16pt}

        What do the ones you circled have in common?\\[18pt]
        \ansline{B and D are measured, so a value can fall anywhere between two others.}\\[18pt]\ansline{A and C are counted --- you get whole numbers and nothing in between.}

\end{enumerate}

\end{document}
